\chapter{Introduction}
\label{chap:introduction}

\section{Purpose}
\label{sec:purpose}

This Software Requirements Specification (SRS) document describes the functional and
non-functional requirements for the \textbf{PQC CBOM Scanner} — a web-based
cybersecurity assessment tool that audits public-facing servers for post-quantum
cryptography (PQC) readiness.

The document is intended for the following audiences:

\begin{itemize}[noitemsep]
  \item \textbf{Development Team} — as a binding technical specification during
        implementation and testing.
  \item \textbf{Security Evaluators} — as a baseline for verifying that the system
        correctly classifies TLS configurations against NIST PQC standards.
  \item \textbf{Hackathon Judges / Reviewers} — as a structured description of the
        project scope, design decisions, and requirements.
  \item \textbf{End Users} — as a reference for understanding system capabilities
        and expected behaviour.
\end{itemize}

\section{Scope}
\label{sec:scope}

The \textbf{PQC CBOM Scanner} (project codename: \textit{quantumCheck}) is a
full-stack cybersecurity tool that:

\begin{enumerate}[noitemsep]
  \item Establishes TLS handshakes with one or more user-supplied hostnames.
  \item Extracts a complete \textbf{Cryptographic Bill of Materials (CBOM)} from
        each TLS session, including cipher suite, key exchange mechanism,
        certificate details, and signature algorithms.
  \item Classifies each scanned asset against a four-tier quantum readiness
        framework aligned to NIST FIPS 203, 204, and 205.
  \item Generates a machine-readable JSON CBOM and a self-contained interactive
        HTML dashboard report.
  \item Optionally discovers and scans subdomains of a root domain using the
        Subfinder tool.
\end{enumerate}

The system is delivered as a \textbf{Dockerised web application} (Flask / Gunicorn)
with a dual interface: a browser-based GUI and a command-line interface.

\textbf{Out of scope:}
\begin{itemize}[noitemsep]
  \item Active exploitation or penetration testing of target hosts.
  \item Certificate Authority (CA) management or certificate issuance.
  \item Remediation automation (the tool provides guidance, not automated fixes).
  \item Support for non-TLS protocols (SSH, IPsec, etc.) in v1.0.
\end{itemize}

\section{Definitions, Acronyms, and Abbreviations}
\label{sec:definitions}

\begin{longtable}{L{3cm} L{11cm}}
  \toprule
  \textbf{Term} & \textbf{Definition} \\
  \midrule
  \endhead
  \bottomrule
  \endfoot
  %
  CBOM     & \textit{Cryptographic Bill of Materials} — a structured inventory
             of all cryptographic primitives used in a system or communication
             session. \\
  PQC      & \textit{Post-Quantum Cryptography} — cryptographic algorithms
             designed to resist attacks from quantum computers. \\
  TLS      & \textit{Transport Layer Security} — the standard protocol for
             encrypted communication over a network. \\
  NIST     & \textit{National Institute of Standards and Technology} — the US
             federal body responsible for cryptographic standards. \\
  FIPS     & \textit{Federal Information Processing Standard} — e.g.\
             FIPS~203 (ML-KEM), FIPS~204 (ML-DSA), FIPS~205 (SLH-DSA). \\
  ML-KEM   & \textit{Module-Lattice-based Key Encapsulation Mechanism}
             (Kyber) — a NIST-standardised PQC key exchange algorithm. \\
  ML-DSA   & \textit{Module-Lattice-based Digital Signature Algorithm}
             (Dilithium) — a NIST-standardised PQC signature scheme. \\
  SLH-DSA  & \textit{Stateless Hash-Based Digital Signature Algorithm}
             (SPHINCS+) — a NIST-standardised PQC signature scheme. \\
  FN-DSA   & \textit{Falcon Digital Signature Algorithm} — another NIST
             PQC signature standard. \\
  ECDSA    & \textit{Elliptic Curve Digital Signature Algorithm}. \\
  ECDH     & \textit{Elliptic Curve Diffie-Hellman} key exchange. \\
  RSA      & \textit{Rivest–Shamir–Adleman} public-key cryptosystem. \\
  HNDL     & \textit{Harvest Now, Decrypt Later} — a threat model where
             adversaries collect encrypted traffic today to decrypt it once
             quantum computers mature. \\
  FQS      & \textit{Fully Quantum Safe} — highest readiness tier. \\
  SSE      & \textit{Server-Sent Events} — a web API for streaming data from
             server to client. \\
  GUI      & \textit{Graphical User Interface}. \\
  CLI      & \textit{Command-Line Interface}. \\
  WSGI     & \textit{Web Server Gateway Interface} — Python web server standard. \\
  SRS      & \textit{Software Requirements Specification}. \\
\end{longtable}

\section{References}
\label{sec:references}

\begin{enumerate}[noitemsep, label={[\arabic*]}]
  \item NIST FIPS 203 — \textit{Module-Lattice-Based Key-Encapsulation Mechanism
        Standard (ML-KEM)}, August 2024.
  \item NIST FIPS 204 — \textit{Module-Lattice-Based Digital Signature Standard
        (ML-DSA)}, August 2024.
  \item NIST FIPS 205 — \textit{Stateless Hash-Based Digital Signature Standard
        (SLH-DSA)}, August 2024.
  \item IEEE Std 830-1998 — \textit{IEEE Recommended Practice for Software
        Requirements Specifications}.
  \item RFC 8446 — \textit{The Transport Layer Security (TLS) Protocol Version 1.3},
        IETF, 2018.
  \item OWASP Cryptographic Failures — \textit{OWASP Top 10 2021, A02}.
  \item ProjectDiscovery Subfinder — \url{https://github.com/projectdiscovery/subfinder}.
  \item NIST NCCoE Migration to Post-Quantum Cryptography Project, SP 1800-38.
\end{enumerate}

\section{Document Overview}
\label{sec:doc_overview}

The remainder of this SRS is structured as follows:

\begin{description}[noitemsep, leftmargin=2cm]
  \item[Chapter~\ref{chap:overall_desc}] provides a product overview, user classes,
        operating environment, and assumptions.
  \item[Chapter~\ref{chap:functional}] specifies all functional requirements grouped
        by system feature.
  \item[Chapter~\ref{chap:external_interfaces}] describes user, hardware, software,
        and communication interfaces.
  \item[Chapter~\ref{chap:nonfunctional}] defines non-functional requirements
        (performance, security, reliability, portability).
  \item[Chapter~\ref{chap:architecture}] presents the system architecture, data
        flow, and deployment model.
  \item[Chapter~\ref{chap:constraints}] lists design constraints, legal
        considerations, and dependencies.
  \item[Chapter~\ref{chap:appendix}] contains the CBOM JSON schema, quantum
        classification logic table, and version history.
\end{description}
